\documentclass[a4paper,11pt]{article}
\usepackage[a4paper,margin=2cm]{geometry}
\usepackage[T1]{fontenc}
\usepackage[utf8]{inputenc}
\usepackage[ngerman,english]{babel}
\usepackage{lmodern,microtype,xcolor,parskip,enumitem}
\usepackage[most]{tcolorbox}
\usepackage{titlesec}
\usepackage[hidelinks]{hyperref}
\definecolor{HeaderBlueDark}{RGB}{70,110,170}
\definecolor{RowBlueLight}{RGB}{235,243,255}
\definecolor{RowBlueDark}{RGB}{220,233,250}
\definecolor{SoftGray}{RGB}{90,90,90}
\setlength{\parindent}{0pt}
\setlist[itemize]{leftmargin=1.4em,itemsep=0.35em,topsep=0.35em}
\titleformat{\section}[block]{\Large\bfseries\color{HeaderBlueDark}}{}{0pt}{}
\newtcolorbox{exbox}{enhanced,breakable,colback=RowBlueLight,colframe=RowBlueDark,boxrule=0.4pt,arc=1.5mm,left=1.2mm,right=1.2mm,top=1mm,bottom=1mm,title={Beispiele \textnormal{(Examples)}},fonttitle=\bfseries,colbacktitle=RowBlueDark,coltitle=black}
\NewDocumentEnvironment{grammarentry}{mm+b}{\begin{tcolorbox}[enhanced,breakable,colback=white,colframe=HeaderBlueDark,title={#1\hfill\normalfont\footnotesize #2},colbacktitle=HeaderBlueDark,coltitle=white,fonttitle=\bfseries]#3\end{tcolorbox}}{}
\newcommand{\GermanText}[1]{\textbf{Deutsch: }#1\par}
\newcommand{\EnglishText}[1]{{\small\color{SoftGray}\textbf{English: }#1\par}}
\newcommand{\ExampleItem}[2]{\item \textbf{#1}\\{\small\color{SoftGray}#2}}
\title{Die Modalpartikel \glqq eben\grqq\\{\large\textnormal{The modal particle eben}}}
\date{}
\begin{document}
\maketitle
\tableofcontents
\section{Bedeutung und Wirkung \textnormal{(Meaning and effect)}}
\begin{grammarentry}{eben}{Modalpartikel / Abtönungspartikel}
\GermanText{Eben präsentiert etwas als folgerichtig, selbstverständlich oder unvermeidlich.}
\EnglishText{Eben presents something as a natural conclusion, self-evident or unavoidable.}
\end{grammarentry}
\section{Grammatische Rolle und Satzarten \textnormal{(Grammar and sentence moods)}}
\begin{grammarentry}{Verwendung}{Satzart und Kontext}
\GermanText{Häufig in Aussagen, Schlussfolgerungen und resignierten Ratschlägen.}
\EnglishText{Common in statements, conclusions and resigned advice.}
\end{grammarentry}
\begin{grammarentry}{Form und Stellung}{Syntax}
\GermanText{Als Modalpartikel ist eben unflektierbar, fordert keinen Kasus und steht normalerweise im Mittelfeld. Allein im Vorfeld eines normalen Aussagesatzes kann eine Modalpartikel nicht stehen. Die Betonung und der Kontext beeinflussen die Wirkung.}
\EnglishText{As a modal particle, eben is invariable, does not govern case and normally appears in the middle field. It cannot occupy the prefield of a normal statement on its own. Stress and context affect its pragmatic force.}
\end{grammarentry}
\section{Typische Beispiele \textnormal{(Typical examples)}}
\begin{exbox}\begin{itemize}
\ExampleItem{Dann musst du eben warten.}{Then you'll just have to wait.}
\ExampleItem{Das ist eben so.}{That's simply the way it is.}
\ExampleItem{Er ist eben noch Anfänger.}{Well, he's still a beginner.}
\ExampleItem{Wenn es regnet, bleiben wir eben hier.}{If it rains, we'll just stay here.}
\end{itemize}\end{exbox}
\section{Abgrenzung zu anderen Verwendungen \textnormal{(Contrast with other uses)}}
\begin{grammarentry}{Wortarten unterscheiden}{Nicht immer Modalpartikel}
\GermanText{Zeitliches eben bedeutet gerade eben; das Adjektiv eben bedeutet flach oder gleichmäßig.}
\EnglishText{Temporal eben means just now; the adjective eben means flat or level.}
\end{grammarentry}
\begin{exbox}\begin{itemize}
\ExampleItem{Ich habe ihn eben gesehen.}{I saw him just now.}
\ExampleItem{Der Boden ist eben.}{The ground is level.}
\end{itemize}\end{exbox}
\section{Lernhinweis \textnormal{(Learning tip)}}
\GermanText{Eine Modalpartikel verändert vor allem die Gesprächsabsicht. Versuchen Sie nicht, jedem Vorkommen eine feste englische Eins-zu-eins-Übersetzung zu geben. Achten Sie auf Satzart, Tonfall und Situation.}
\EnglishText{A modal particle mainly affects the speaker's intention. Do not try to give it one fixed English translation in every context. Pay attention to sentence mood, intonation and situation.}
\section{Quellen \textnormal{(Sources)}}
\begin{itemize}
\item IDS Mannheim: \url{https://grammis.ids-mannheim.de/terminologie/2}
\item Duden: \url{https://www.duden.de/sprachwissen/sprachratgeber/Modalpartikeln}
\end{itemize}
\end{document}
